% Shared by both versions: the tree, the tree output and the exit figure.
% Input inside the protocol section (full) or "How LayerTree Works" (short).

\subsection{The Tree}\label{sec:tree}

We write $\mathsf{le}_{k}(x)$ for the $k$-byte little-endian encoding of
$x$. A leaf is a fixed 116-byte record
\begin{multline*}
\lesf(\mathsf{amount}) \,\|\, \mathsf{exit\_key} \,\|\, \mathsf{owner\_key}
\,\|\, \mathsf{evictor\_key} \\
\,\|\, \mathsf{le}_{32}(\mathsf{evict\_after})
\,\|\, \mathsf{le}_{32}(\mathsf{nonce}) \,\|\, \mathsf{le}_{32}(\mathsf{seq}),
\end{multline*}
with keys as 32-byte x-only points, $\mathsf{evict\_after}$ a block
height, $\mathsf{nonce}$ a per-leaf identifier and $\mathsf{seq}$ a
counter advanced by vouchers (\iffullversion Appendix~\ref{app:claim}\else \S\ref{sec:vouchers}\fi). A \emph{node} is
a pair $(h, s)$ of a 32-byte hash and an unsigned sum. A position holding
leaf $\ell$ has node $(\Hleaf(\ell), \ell.\mathsf{amount})$; an empty
position has node $\Emp_0 = (0^{32}, 0)$. The parent of $(h_L, s_L)$ and
$(h_R, s_R)$ is
\[
\bigl(\Hnode(h_L \,\|\, \encs(s_L) \,\|\, h_R \,\|\, \encs(s_R)),\; s_L + s_R\bigr),
\qquad \encs(s) = \lesf(s) \,\|\, \mathtt{0x01}.
\]
Every hash is a BIP-340 tagged hash with its own tag; $\encs(s)$ is the
minimal BIP-441 encoding of $s + 2^{64}$, produced from a computed sum
with one addition\iffullversion{} (Lemma~\ref{lem:enc})\fi. $\Emp_h$ is the node of the
empty tree of height $h$.

The \emph{sibling path} $\pi$ of a position is the list of $D$ pairs
$(\mathsf{dir}_k, (h_k, s_k))$, one per level, giving the side the
position lies on and the sibling node not taken. $\fold(n, \pi)$ combines
a node with a path from the bottom up, hashing with each sibling in the
order the direction gives and adding the sums. The honest path folds a
leaf's node to the root, and folding a \emph{different} node up the
\emph{same} path yields the root with that position replaced
\iffullversion(Theorem~\ref{thm:binding})\else(\S\ref{sec:proved})\fi.

\paragraph{Profiles.}
\iffullversion
In the \emph{dense} profile positions are assigned by the operator and a
path carries explicit direction bits; each level costs 88 bytes of script
and 41 bytes of witness. In the \emph{sparse} profile a leaf's position is
the low $D$ bits of $\Hpos(\mathsf{owner\_key} \,\|\, \mathsf{nonce})$,
direction bits are derived in script from the position hash, and empty
siblings are elided from the witness (the script substitutes $\Emp_k$),
so witness data grows with $\log_2 N$ rather than with $D$. It is the
Merkle-sum sparse Merkle tree of Taproot Assets~\cite{mssmt} truncated to
$D = 32$ or $64$, at 1.7 to 2.8 times the dense per-exit cost with
$2^{16}$ members (Table~\ref{tab:sizes}). Dense is the default, and the
rest of the paper measures it unless it says otherwise.
\else
In the \emph{dense} profile the operator assigns positions and a path
carries direction bits. In the \emph{sparse} profile a leaf's position
is derived by hashing its owner key and nonce, as in the Merkle-sum
sparse Merkle tree of Taproot Assets~\cite{mssmt}, truncated to
$D = 32$ or $64$, at 1.7 to 2.8 times the dense exit cost. Dense is the
default and is what we measure.
\fi

\subsection{The Tree Output}\label{sec:output}

Let $r$ be the root hash and $S$ the root sum. The tree output is a
taproot output with internal key
$P = \mathsf{lift\_x}\bigl(x(\NUMS + \Hstate(r)\cdot G)\bigr)$, where
$\NUMS$ is the BIP-341 nothing-up-my-sleeve point and $\mathsf{lift\_x}$
the BIP-340 even-$y$ normalisation, so the key path is unspendable and
every spend goes through a script. Its taptree has one leaf per
predicate; its Merkle root $m$ identifies the rule set and is fixed for
the tree's life. A \emph{tree output} for state $(r, S)$ has scriptPubKey
$\TreeSpk(r, m)$ and amount $S$. Every predicate recomputes
$\TreeSpk(r, m)$ from its witness and checks it against the spent input,
binding the root the proofs are checked against and the script set the
successor carries. FULL, MID and MIN predicates read input~0, so two such
trees cannot be spent in one transaction; a CCV tree's input may sit at
any index.

\begin{figure}[tbp]
\centering
\begin{tikzpicture}[
  node distance=4mm and 2mm,
  every node/.style={font=\scriptsize},
  tn/.style={draw, rounded corners=1pt, minimum width=9mm, minimum height=4.5mm, inner sep=1pt},
  sib/.style={tn, fill=black!12},
  path/.style={tn, thick},
  leaf/.style={tn, minimum width=7mm},
  box/.style={draw, minimum height=7mm, align=center, inner sep=2pt},
  arr/.style={-{Latex[length=1.6mm]}, thin}
]
\node[path] (root) {$(r,S)$};
\node[sib, below left=of root, xshift=-8mm] (n0) {$(h_0,s_0)$};
\node[path, below right=of root, xshift=8mm] (n1) {$(h_1,s_1)$};
\node[leaf, below left=of n0, xshift=2mm] (l00) {$\ell_0$};
\node[leaf, below right=of n0, xshift=-2mm] (l01) {$\ell_1$};
\node[path, below left=of n1, xshift=2mm] (l10) {$\ell_2$};
\node[sib, below right=of n1, xshift=-2mm] (l11) {$\Emp_0$};
\draw (root) -- (n0); \draw (root) -- (n1);
\draw (n0) -- (l00); \draw (n0) -- (l01);
\draw (n1) -- (l10); \draw (n1) -- (l11);
\node[below=1.5mm of l01, xshift=8mm, align=center] {shaded: sibling path of $\ell_2$,\quad $\pi = [(\mathsf{R}, (h_0,s_0)),\ (\mathsf{L}, \Emp_0)]$};
\begin{scope}[xshift=37mm, yshift=1mm]
\node[box, minimum width=24mm] (in0) {input 0\\ $\TreeSpk(r,m)$, $S$};
\node[box, minimum width=24mm, below=2.5mm of in0] (in1) {input 1 (fee)};
\node[box, minimum width=26mm, right=5mm of in0] (out0) {output 0\\ $\TreeSpk(r',m)$, $S-a$};
\node[box, minimum width=26mm, below=2.5mm of out0] (out1) {output 1\\ $\mathrm{P2TR}(\ell_2.\mathsf{exit\_key})$, $a$};
\node[box, minimum width=26mm, minimum height=5mm, below=2.5mm of out1] (out2) {change, anchors (free)};
\draw[arr] (in0) -- (out0);
\draw[arr] (in0) -- (out1);
\draw[arr] (in1) -- (out2);
\node[above=0.5mm of in0, xshift=17mm, font=\scriptsize\itshape] {\tr{EXIT} transaction for $\ell_2$, amount $a$};
\end{scope}
\end{tikzpicture}
\caption{A depth-2 tree and the exit of $\ell_2$. The predicate checks $\fold((\Hleaf(\ell_2),a),\pi) = (r,S)$, computes $r' = \fold(\Emp_0,\pi)$, verifies a signature by $\ell_2.\mathsf{owner\_key}$, and constrains outputs 0 and 1. Everything else in the transaction is free.}
\label{fig:exit}
\end{figure}
